\documentclass[UTF8,fontset=fandol,openany,zihao=-4]{ctexbook}
% Shared lecture-note preamble for Big Data and Management Decision-Making.
% Chapter files follow latex_config.md and remain independently compilable.

\usepackage[a4paper,margin=2.6cm]{geometry}
\usepackage{amsmath,amssymb,amsthm,mathtools,bm}
\usepackage{xcolor}
\usepackage[most]{tcolorbox}
\usepackage{booktabs,longtable,array}
\usepackage{enumitem}
\usepackage{hyperref}
\usepackage{listings}
\usepackage{caption}
\usepackage{tikz}
\usetikzlibrary{arrows.meta,positioning,shapes.geometric}

\newcommand{\BDMFandolPath}{/usr/local/texlive/2025/texmf-dist/fonts/opentype/public/fandol/}
\setCJKmainfont[
  Path=\BDMFandolPath,
  UprightFont=FandolSong-Regular.otf,
  BoldFont=FandolSong-Regular.otf,
  ItalicFont=FandolKai-Regular.otf,
  BoldItalicFont=FandolKai-Regular.otf
]{FandolSong-Regular.otf}
\setCJKsansfont[
  Path=\BDMFandolPath,
  UprightFont=FandolSong-Regular.otf,
  BoldFont=FandolSong-Regular.otf
]{FandolSong-Regular.otf}
\setCJKmonofont[
  Path=\BDMFandolPath,
  UprightFont=FandolSong-Regular.otf,
  BoldFont=FandolSong-Regular.otf
]{FandolSong-Regular.otf}
\setmonofont{Menlo}
\linespread{1.13}

\hypersetup{
  colorlinks=true,
  linkcolor=blue!60!black,
  citecolor=blue!60!black,
  urlcolor=blue!60!black,
  pdfauthor={大数据与管理决策基础},
  pdfsubject={大数据与管理决策基础课程讲义}
}

\ctexset{
  chapter = {
    format = \huge\mdseries,
    name = {第,讲},
    number = \arabic{chapter},
    aftername = \quad
  },
  section = {format = \Large\mdseries},
  subsection = {format = \large\mdseries},
  subsubsection = {format = \normalsize\mdseries}
}

\setlist[itemize]{leftmargin=2em,itemsep=0.25em,topsep=0.25em}
\setlist[enumerate]{leftmargin=2em,itemsep=0.25em,topsep=0.25em}

\definecolor{BDMBlue}{RGB}{22,44,73}
\definecolor{BDMTeal}{RGB}{22,119,119}
\definecolor{BDMGray}{RGB}{76,84,95}
\definecolor{BDMLight}{RGB}{245,247,250}
\definecolor{BDMAmber}{RGB}{181,111,35}
\definecolor{CodeBg}{RGB}{248,248,248}

\lstdefinestyle{pythonstyle}{
  basicstyle=\ttfamily\small,
  backgroundcolor=\color{CodeBg},
  keywordstyle=\color{BDMBlue},
  commentstyle=\color{BDMGray},
  stringstyle=\color{BDMAmber},
  showstringspaces=false,
  breaklines=true,
  frame=single,
  rulecolor=\color{black!20},
  columns=fullflexible
}

\lstdefinestyle{sqlstyle}{
  language=SQL,
  basicstyle=\ttfamily\small,
  backgroundcolor=\color{CodeBg},
  keywordstyle=\color{BDMBlue},
  commentstyle=\color{BDMGray},
  stringstyle=\color{BDMAmber},
  showstringspaces=false,
  breaklines=true,
  frame=single,
  rulecolor=\color{black!20},
  columns=fullflexible
}

\lstset{style=pythonstyle}
\renewcommand{\lstlistingname}{代码}
\renewcommand{\bibname}{参考文献}

\theoremstyle{definition}
\newtheorem{definition}{定义}[chapter]
\newtheorem{example}[definition]{例}
\newtheorem{exercise}[definition]{习题}
\newtheorem{convention}[definition]{约定}

\theoremstyle{remark}
\newtheorem{remark}[definition]{评注}

\theoremstyle{plain}
\newtheorem{theorem}[definition]{定理}
\newtheorem{proposition}[definition]{命题}
\newtheorem{lemma}[definition]{引理}
\newtheorem{corollary}[definition]{推论}

\tcolorboxenvironment{definition}{
  colback=blue!2!white,
  colframe=blue!45!black,
  boxrule=0.4pt,
  arc=1.2mm,
  breakable
}

\tcolorboxenvironment{theorem}{
  colback=orange!3!white,
  colframe=orange!55!black,
  boxrule=0.4pt,
  arc=1.2mm,
  breakable
}

\tcolorboxenvironment{proposition}{
  colback=orange!2!white,
  colframe=orange!50!black,
  boxrule=0.4pt,
  arc=1.2mm,
  breakable
}

\tcolorboxenvironment{lemma}{
  colback=orange!2!white,
  colframe=orange!45!black,
  boxrule=0.4pt,
  arc=1.2mm,
  breakable
}

\tcolorboxenvironment{corollary}{
  colback=orange!2!white,
  colframe=orange!45!black,
  boxrule=0.4pt,
  arc=1.2mm,
  breakable
}

\newtcolorbox{chapterbox}{
  colback=gray!5!white,
  colframe=gray!55!black,
  boxrule=0.4pt,
  arc=1.2mm,
  breakable
}

\newtcolorbox{modernbox}[1][应用接口]{
  colback=green!3!white,
  colframe=green!45!black,
  boxrule=0.4pt,
  arc=1.2mm,
  title={#1},
  fonttitle=\mdseries,
  breakable
}

\newcommand{\R}{\mathbb R}
\newcommand{\N}{\mathbb N}
\newcommand{\E}{\mathbb E}
\newcommand{\Prob}{\mathbb P}
\newcommand{\dd}{\,\mathrm d}
\newcommand{\eps}{\varepsilon}
\newcommand{\abs}[1]{\left\lvert #1\right\rvert}
\newcommand{\norm}[1]{\left\lVert #1\right\rVert}
\newcommand{\argmin}{\operatorname*{arg\,min}}
\newcommand{\argmax}{\operatorname*{arg\,max}}


\title{《大数据与管理决策基础》\\[0.5em]\large 第6讲：统计推断、A/B 测试与管理判断}
\author{课程讲义}
\date{}

\begin{document}

\maketitle
\tableofcontents

\setcounter{chapter}{5}
\chapter{统计推断、A/B 测试与管理判断}
\label{ch:week06-ab-testing}

\begin{chapterbox}
第5讲已经说明，dashboard 能把经营状态组织为可比较、可下钻、可行动的信息界面。但管理者看到差异之后还必须继续追问：这种差异是可重复的经营改进，还是有限样本中的随机波动？本讲把描述统计推进到统计推断和随机实验。讲义首先说明总体、样本、估计量、标准误和置信区间的含义；随后给出 A/B 测试的基本结构，包括随机分配、实验单元、处理、对照、主指标、护栏指标和最小可检测效应；再用差异估计量、功效分析、重复窥探、多重比较和实验报告治理说明如何从实验结果形成审慎的管理判断；最后通过结账页改版案例展示代码、报告和决策复盘的完整接口。
\end{chapterbox}

\begin{modernbox}[课程接口]
本讲把 dashboard 中的指标差异推进到可证据化的行动判断：若第5讲回答“经营状态看起来如何”，第6讲回答“观察到的差异是否足以支持改变”。它也是第8讲因果推断的前置章节，因为随机实验提供了最清晰的因果识别基准。
\end{modernbox}

\section{学习目标}

完成本讲后，学生应能够：
\begin{enumerate}
  \item 区分描述统计、统计推断和因果评估的任务边界；
  \item 解释总体、样本、估计量、标准误、置信区间、p 值、显著性水平和检验功效；
  \item 设计一个面向管理行动的 A/B 测试，包括实验单元、处理、对照、随机化机制、主指标和护栏指标；
  \item 使用差异估计量评估转化率、收入、时延和客服工单等指标的实验效果；
  \item 理解最小可检测效应、样本量不足、重复窥探和多重指标带来的决策风险；
  \item 识别实验污染、样本比例失衡、指标口径漂移和选择性报告等常见问题；
  \item 将实验结果写成包含业务解释、统计不确定性和行动建议的管理报告。
\end{enumerate}

\section{从 Dashboard 差异到统计推断}

第5讲的 dashboard 让管理者看见了差异。例如某一天收入升高，某个地区准时率下降，某个新版页面转化率似乎更好。差异本身不是结论。管理者真正需要判断的是：若在相同业务环境下重复观察或重复实施行动，这种差异是否仍然存在？若把观察到的差异用于推广、投放、流程改版或资源配置，组织是否有足够证据承担相应成本和风险？

描述统计面对的是已经观察到的数据。统计推断面对的是未完全观察的总体和未来重复样本。若 \(Y\) 表示某个管理指标，例如是否转化、订单收入或页面加载时间，样本均值
\[
  \bar Y=\frac{1}{n}\sum_{i=1}^n Y_i
\]
只是总体均值 \(\mu=\E[Y]\) 的估计。不同样本会产生不同的 \(\bar Y\)。因此，管理报告不能只写“B 组转化率比 A 组高 20 个百分点”，还要说明这种估计的不确定性、样本量、适用对象和业务阈值。

\begin{definition}[统计推断]
统计推断是利用样本数据对未完全观察的总体参数或未来重复样本行为作出带不确定性的判断。它关注估计量、标准误、置信区间、检验和决策风险。
\end{definition}

统计推断和因果评估也不同。统计推断可以描述样本差异的不确定性；因果评估还需要说明差异是否由某个行动造成。A/B 测试之所以重要，是因为随机化同时服务于两件事：减少选择偏差，并使差异估计量具有可分析的抽样不确定性。

\section{总体、样本、估计量与标准误}

总体是管理者关心但无法完全观察的对象集合。例如未来所有访问结账页的访客、本季度所有潜在订单、未来 30 天可能接触新推荐策略的用户。样本是已经观察到的那部分对象。参数是总体中的未知量，估计量是从样本计算出来的统计量。

\begin{definition}[标准误]
设 \(\hat\theta\) 是总体参数 \(\theta\) 的估计量。标准误是估计量在重复抽样下的标准差，记为
\[
  \operatorname{SE}(\hat\theta)=\sqrt{\operatorname{Var}(\hat\theta)}.
\]
在实际数据中，标准误通常需要用样本方差或模型估计来近似。
\end{definition}

对均值而言，若样本方差为 \(s^2\)，则
\[
  \widehat{\operatorname{SE}}(\bar Y)=\frac{s}{\sqrt n}.
\]
这个公式给出一个重要管理含义：样本量增加会降低不确定性，但速度是 \(1/\sqrt n\)，不是 \(1/n\)。若希望把标准误减半，通常需要约四倍样本量。样本量不足时，dashboard 上看起来很大的差异也可能不足以支持推广。

\section{随机实验的管理语言}

A/B 测试是一种随机对照实验。它把符合条件的实验单元随机分配到对照组和处理组，然后比较两组结果。随机化的作用不是让两组在每个变量上完全相同，而是让处理分配与潜在结果在设计上独立。这样，组间平均差异可以解释为处理效果的估计，而不是选择偏差的直接产物。

\begin{definition}[A/B 测试的基本元素]
一个面向管理决策的 A/B 测试至少包含：
\begin{enumerate}
  \item 实验单元：被随机分配的对象，如访客、订单、门店、销售人员或客户；
  \item 处理与对照：\(T_i=1\) 表示处理组，\(T_i=0\) 表示对照组；
  \item 主指标：用于判断实验是否达到核心目标的指标；
  \item 护栏指标：用于防止局部优化损害系统安全、体验或长期价值的指标；
  \item 分析窗口：从分配到观察结果的时间范围；
  \item 纳入和排除规则：哪些对象有资格进入实验，哪些记录应被排除；
  \item 决策规则：何时推广、继续实验、停止实验或回滚。
\end{enumerate}
\end{definition}

在企业场景中，A/B 测试经常用于页面改版、推荐策略、优惠规则、履约承诺、客服话术和运营流程。好的实验设计不是“把新版给一半用户看看”，而是把管理问题写成可执行的因果问题。例如：
\[
  \text{是否应将新版结账页推广给全部访客？}
\]
可以被转化为：
\[
  \text{新版结账页是否提高访客转化率，并且不显著恶化页面加载时间和客服工单率？}
\]

\section{潜在结果与随机化}

为了说明 A/B 测试为什么能够支持因果判断，引入潜在结果记号。对实验单元 \(i\)，令 \(Y_i(1)\) 表示接受处理时的结果，\(Y_i(0)\) 表示接受对照时的结果。个体处理效应为
\[
  \tau_i=Y_i(1)-Y_i(0).
\]
但同一个实验单元不能同时处于处理和对照状态，所以 \(\tau_i\) 无法被直接观察。实验目标通常是平均处理效应：
\[
  \tau=\E[Y(1)-Y(0)].
\]

若处理分配 \(T_i\) 是随机的，且实验执行没有系统性污染，那么处理组均值与对照组均值之差
\[
  \hat\tau
  =\bar Y_1-\bar Y_0
  =\frac{1}{n_1}\sum_{i:T_i=1}Y_i
   -\frac{1}{n_0}\sum_{i:T_i=0}Y_i
\]
就是 \(\tau\) 的自然估计量。

\begin{proposition}[随机化差异估计量的无偏性]
若处理分配与潜在结果 \((Y_i(1),Y_i(0))\) 独立，且处理组与对照组均有正概率被观察，则
\[
  \E[\bar Y_1-\bar Y_0]=\E[Y(1)-Y(0)].
\]
\end{proposition}

这个命题说明了实验设计优先于复杂模型的原因。若随机化设计被破坏，例如高价值客户更容易看到新版，或者某些渠道只进入处理组，那么组间差异就同时包含处理效果和选择偏差。此时即使后续使用复杂模型，也不能轻易恢复实验本来应有的解释力。

\section{均值差异、置信区间与 p 值}

对连续指标或二元指标，都可以从均值差异开始。二元转化指标 \(Y_i\in\{0,1\}\) 的均值就是转化率。设处理组和对照组样本均值分别为 \(\bar Y_1,\bar Y_0\)，样本方差分别为 \(s_1^2,s_0^2\)，则差异估计量的标准误可写为
\[
  \widehat{\operatorname{SE}}(\hat\tau)
  =
  \sqrt{\frac{s_1^2}{n_1}+\frac{s_0^2}{n_0}}.
\]
在大样本近似下，\(100(1-\alpha)\%\) 置信区间为
\[
  \hat\tau
  \pm
  z_{1-\alpha/2}\widehat{\operatorname{SE}}(\hat\tau).
\]

\begin{definition}[p 值]
在给定原假设 \(H_0\) 的条件下，p 值是观察到当前统计量或更极端统计量的概率。它不是原假设为真的概率，也不是处理有效的概率。
\end{definition}

本课程要求学生谨慎使用 p 值。p 值可以作为证据强弱的一个摘要，但它不能替代业务成本、实施风险、护栏指标、样本代表性和长期效果判断。尤其在管理场景中，“统计显著”不等于“值得推广”，“不显著”也不等于“没有效果”。小样本实验可能看起来方向正确，但由于不确定性较大，仍需要继续实验或扩大样本。

\section{显著性、功效与最小可检测效应}

实验设计不能等结果出来后才讨论样本量。若样本量太小，实验可能长期处于“看起来不错但无法判断”的状态。显著性水平 \(\alpha\) 控制第一类错误，即在没有真实效果时错误推广的概率。功效 \(1-\beta\) 控制在真实效果达到某个幅度时能够发现它的概率。二者共同决定实验对风险的态度。

最小可检测效应，通常记为 MDE，是组织希望实验能够可靠识别的最小业务差异。例如，若当前转化率为 35\%，管理者可能认为至少提升 10 个百分点才值得承担改版风险。

对两组比例差异，一个常用的每组样本量近似为
\[
n
\approx
\frac{
\left[
z_{1-\alpha/2}\sqrt{2\bar p(1-\bar p)}
+
z_{1-\beta}\sqrt{p_0(1-p_0)+p_1(1-p_1)}
\right]^2
}{
(p_1-p_0)^2
},
\]
其中 \(p_0\) 是基准转化率，\(p_1=p_0+\delta\)，\(\delta\) 是 MDE，\(\bar p=(p_0+p_1)/2\)，\(1-\beta\) 是检验功效。

\begin{modernbox}[管理解释]
MDE 不是统计软件的参数，而是管理层对“多大改善值得行动”的声明。若 MDE 设得过小，所需样本量可能远超业务可承受范围；若 MDE 设得过大，实验可能错过虽小但长期价值很高的改进。
\end{modernbox}

\section{护栏指标与实验决策}

许多失败的实验并不是主指标没有提升，而是主指标提升的同时损害了系统其他部分。例如新版结账页提高了转化率，却让页面加载时间显著变慢；推荐策略提高了点击率，却降低了复购；优惠规则提高了订单数，却压缩了毛利。

\begin{definition}[护栏指标]
护栏指标是用于约束实验推广的风险指标。它不一定是实验的主要优化目标，但一旦明显恶化，实验结论就不能只依据主指标做推广决策。
\end{definition}

护栏指标常见于以下几类：
\begin{longtable}{>{\raggedright\arraybackslash}p{0.20\linewidth}>{\raggedright\arraybackslash}p{0.34\linewidth}>{\raggedright\arraybackslash}p{0.34\linewidth}}
\caption{A/B 测试中的常见护栏指标}\\
\toprule
类型 & 指标例子 & 管理含义\\
\midrule
\endfirsthead
\toprule
类型 & 指标例子 & 管理含义\\
\midrule
\endhead
体验护栏 & 页面加载时间、投诉率、客服工单率 & 避免短期转化提升损害用户体验。\\
财务护栏 & 毛利率、补贴成本、退货率 & 避免收入提升来自不可持续补贴或高退货。\\
运营护栏 & 履约延迟、库存缺货、人工处理量 & 避免前端策略把压力转移到供应链或运营团队。\\
长期护栏 & 复购率、留存率、订阅取消率 & 避免短期行为伤害长期客户价值。\\
\bottomrule
\end{longtable}

因此，实验报告至少应同时呈现主指标和护栏指标的方向、置信区间、业务阈值和样本量。若主指标置信区间跨越 0，或护栏指标出现明显恶化，稳妥的管理建议通常不是立即全量推广，而是继续实验、缩小适用范围、优化方案或进行分层分析。

\section{实验质量检查}

一个实验即使能计算 p 值，也未必是可信实验。质量检查应在分析前完成，至少包括以下内容。

\begin{longtable}{>{\raggedright\arraybackslash}p{0.24\linewidth}>{\raggedright\arraybackslash}p{0.34\linewidth}>{\raggedright\arraybackslash}p{0.30\linewidth}}
\caption{实验质量检查}\\
\toprule
检查项 & 问题 & 风险\\
\midrule
\endfirsthead
\toprule
检查项 & 问题 & 风险\\
\midrule
\endhead
随机分配 & 处理组和对照组是否按预设比例进入？ & 样本比例失衡可能意味着埋点或分流错误。\\
实验单元 & 分配单元和分析单元是否一致？ & 同一用户跨组会污染处理效果。\\
数据完整性 & 指标事件是否缺失或重复？ & 差异可能来自记录错误。\\
分析窗口 & 是否所有单元都有相同观察窗口？ & 后进入实验的单元可能被低估结果。\\
口径冻结 & 主指标和护栏指标是否实验前定义？ & 事后挑选指标会放大偶然显著。\\
\bottomrule
\end{longtable}

样本比例失衡常被称为 SRM。若原计划 50/50 分流，但实际进入 B 组的样本显著偏少，可能说明随机化、日志采集或资格过滤出了问题。此时应先排查实验机制，而不是直接解释组间差异。

\section{重复窥探与多重比较}

A/B 测试在企业中常被频繁查看。每天打开实验看一次结果，若显著就停止，看似谨慎，实际会提高误报率。原因是每次查看都相当于给偶然波动一次“被选中”的机会。除非使用序贯检验或预先规定的中期分析规则，否则应在实验开始前设定分析窗口和停止规则。

多重比较也会带来类似问题。若同时查看 20 个指标、10 个分群和多个时间窗口，即使没有真实效果，也很可能出现若干“显著”结果。课程项目中不要求学生掌握复杂校正方法，但必须在报告中说明哪些指标是预先定义的主指标，哪些是探索性分析。

\begin{remark}
探索性分群并非禁止。它的正确定位是提出后续假设，而不是直接作为推广依据。若发现某个客户分层效果很强，应设计新的确认性实验或在后续样本中验证。
\end{remark}

\section{实验报告与决策 Memo}

一个可复核的实验报告不应只是统计输出。最低结构如下：
\begin{enumerate}
  \item 业务问题：本实验要支持什么管理决策；
  \item 实验设计：实验单元、处理、对照、随机化、纳入规则和分析窗口；
  \item 指标体系：主指标、护栏指标、MDE、显著性水平和功效；
  \item 数据质量：样本量、分流比例、缺失、重复、异常和口径版本；
  \item 统计结果：效应估计、标准误、置信区间、p 值和样本量解释；
  \item 管理解释：收益、成本、风险、适用范围和不确定性；
  \item 行动建议：推广、继续实验、回滚、改版或分层复验；
  \item 复盘入口：实验代码、数据批次、报告版本和最终决策。
\end{enumerate}

这份 memo 的作用是让团队在数周或数月后仍能理解当时为什么做出某个决定。实验治理的目标不是追求每次实验都显著，而是让每次行动都有可复核的证据链。

\section{第6周实验案例}

本周样例数据位于：
\begin{center}
\path{data/sample/ab_test_checkout.csv}
\end{center}
该数据模拟一个结账页 A/B 测试。实验单元为访客，A 组为旧版结账页，B 组为新版结账页。主指标是是否转化，辅助指标是每访客收入，护栏指标包括页面加载时间和客服工单率。

配套代码位于：
\begin{center}
\path{src/bdm_decision/cases/week06_ab_testing.py}
\end{center}

最小运行方式为：
\begin{lstlisting}[style=pythonstyle,caption={运行第6周 A/B 测试案例}]
PYTHONPATH=src \
python3 src/bdm_decision/cases/week06_ab_testing.py
\end{lstlisting}

核心结果如下：
\begin{itemize}
  \item A 组样本量为 20，转化率为 0.3500，每访客收入为 16.95；
  \item B 组样本量为 20，转化率为 0.5500，每访客收入为 31.65；
  \item 转化率差异为 0.2000，p 值为 0.2059，95\% 置信区间跨越 0；
  \item 每访客收入差异为 14.70，95\% 置信区间跨越 0；
  \item 页面加载时间平均增加 51.65 毫秒，护栏指标显示处理组可能恶化体验；
  \item 若基准转化率为 0.35，希望检测 10 个百分点的 MDE，在 80\% 功效下每组约需 376 个样本。
\end{itemize}

这个案例的教学目的不是证明 B 组无效，而是说明管理决策应同时考虑方向、幅度、不确定性、样本量和护栏风险。在当前小样本下，合理建议是继续实验或扩大样本，而不是立刻全量推广。

\section*{本讲小结}
\addcontentsline{toc}{section}{本讲小结}

本讲将 dashboard 中的指标差异转化为统计推断问题。随机实验通过设计上的随机化减少选择偏差，差异估计量提供处理效果的自然估计，标准误和置信区间刻画估计不确定性，p 值提供原假设下的极端性摘要，功效和 MDE 帮助实验在开始前形成可执行计划。管理判断不能只看显著性，还必须同时检查主指标、护栏指标、样本代表性、业务成本、推广风险和实验治理记录。

\section*{习题}
\addcontentsline{toc}{section}{习题}

\begin{exercise}
围绕 \path{data/sample/ab_test_checkout.csv} 完成下列任务：
\begin{enumerate}
  \item 分别计算 A 组和 B 组的样本量、转化率、每访客收入、平均订单价值、页面加载时间和客服工单率；
  \item 使用差异估计量计算 B 组相对 A 组的转化率提升，并给出 95\% 置信区间；
  \item 说明为什么“B 组转化率高 20 个百分点”不足以直接支持全量推广；
  \item 若管理层认为至少提升 5 个百分点才值得推广，请重新估算所需每组样本量；
  \item 检查该实验是否存在样本比例失衡、护栏恶化或指标口径风险；
  \item 为该实验写一段不超过 300 字的管理建议，必须同时提及主指标、护栏指标和下一步行动。
\end{enumerate}
\end{exercise}

\section*{常见误区}
\addcontentsline{toc}{section}{常见误区}

\begin{enumerate}
  \item 把 A/B 测试理解为“上线两个版本看哪个数值高”，忽视随机化、样本量和分析窗口；
  \item 将 p 值误解为处理无效或有效的概率；
  \item 在实验过程中反复窥探结果，并在偶然显著时提前停止；
  \item 只报告转化率，不报告分母、置信区间和护栏指标；
  \item 忽视实验污染和分配单元错误，例如同一用户跨组；
  \item 将短期收入提升解释为长期客户价值提升；
  \item 在没有预先定义主指标和 MDE 的情况下事后挑选有利结果。
\end{enumerate}

\addcontentsline{toc}{chapter}{参考文献}
\begin{thebibliography}{99}
\raggedright
\bibitem{fisher1935}
Fisher, R. A. (1935).
\emph{The Design of Experiments}.
Oliver and Boyd.

\bibitem{neyman1923}
Neyman, J. (1923/1990).
On the application of probability theory to agricultural experiments.
\emph{Statistical Science}, 5(4), 465--480.

\bibitem{cox1958}
Cox, D. R. (1958).
\emph{Planning of Experiments}.
Wiley.

\bibitem{box2005}
Box, G. E. P., Hunter, J. S., and Hunter, W. G. (2005).
\emph{Statistics for Experimenters: Design, Innovation, and Discovery} (2nd ed.).
Wiley.

\bibitem{kohavi2020}
Kohavi, R., Tang, D., and Xu, Y. (2020).
\emph{Trustworthy Online Controlled Experiments: A Practical Guide to A/B Testing}.
Cambridge University Press.

\bibitem{imbens2015}
Imbens, G. W., and Rubin, D. B. (2015).
\emph{Causal Inference for Statistics, Social, and Biomedical Sciences}.
Cambridge University Press.

\bibitem{gerber2012}
Gerber, A. S., and Green, D. P. (2012).
\emph{Field Experiments: Design, Analysis, and Interpretation}.
W. W. Norton.

\bibitem{wasserstein2016}
Wasserstein, R. L., and Lazar, N. A. (2016).
The ASA statement on p-values: Context, process, and purpose.
\emph{The American Statistician}, 70(2), 129--133.
\end{thebibliography}

\end{document}
